%=====================================================================
\chapter{Benchmarking, Simulation and Performance Evaluation}\label{ch:benchmarking}
%=====================================================================
\locator{3}{Part~3 --- Research Designs for Computing}

\begin{outcomes}
By the end of this chapter you will be able to:
\begin{tight}
  \item \textbf{Design} a benchmark whose workload is defensibly representative.
  \item \textbf{Measure} performance with a methodology that survives
    non-determinism.
  \item \textbf{Choose} the correct summary statistic, and explain why the
    arithmetic mean of ratios is wrong.
  \item \textbf{Report} performance results with uncertainty rather than as
    point values.
  \item \textbf{Validate} a simulation model, and distinguish verification from
    validation.
\end{tight}
\end{outcomes}

\begin{prereq}
\chref{experiments} for design and randomisation --- a benchmark run is an
experiment whose subjects are configurations. \chref{questions} for the
comparison element, which benchmarking makes concrete.
\end{prereq}

\begin{keyterms}
benchmark $\bullet$ workload characterisation $\bullet$ representativeness
$\bullet$ warm-up and steady state $\bullet$ non-determinism $\bullet$
repetition and replication $\bullet$ geometric mean $\bullet$ speedup
$\bullet$ confidence interval $\bullet$ verification and validation $\bullet$
transient removal $\bullet$ benchmark overfitting
\end{keyterms}

\section{A Benchmark Is an Argument About Representativeness}

Every performance claim has the same structure: on \emph{this} workload, on
\emph{this} platform, under \emph{this} measurement methodology, X was faster
than Y by \emph{this much}. Each qualifier is a place the claim can fail to
transfer, and the benchmark is the argument that the qualifiers do not matter
as much as one might fear.

\begin{center}
\small
\begin{tabular}{@{}L{3.2cm}L{11.1cm}@{}}
\toprule
\textbf{Question} & \textbf{What a good answer looks like} \\
\midrule
Representative of what?
  & A named target population of real workloads, characterised on dimensions
    that matter --- request mix, data size and skew, concurrency, arrival
    pattern \\
How do you know?
  & Evidence: production traces, published characterisations, a survey of
    practitioners. Not intuition \\
What does it exclude?
  & The workload regimes deliberately not covered, and what would happen there \\
Can it be re-run?
  & Version-pinned software, published configuration, released harness and
    data (\chref{reproducibility}) \\
\bottomrule
\end{tabular}
\end{center}

\begin{pitfall}[Building the benchmark that your system wins]
The characteristic failure, and it is usually unconscious. You develop the
system, testing on workloads that exercise the parts you optimised. Those
workloads become the evaluation. The comparison is real, the measurement is
careful, and the result is uninformative --- because the workload was selected,
implicitly, by which cases the system handles well.

\medskip
Two defences. Fix the workload set \emph{before} building, from an external
characterisation, and write it into the protocol. And include at least one
workload where you expect your approach to lose, then report that it did. A
paper that shows where its own contribution fails is dramatically more credible
than one that wins everywhere, and reviewers know it.
\end{pitfall}

\section{Measurement Under Non-Determinism}

Modern systems are not deterministic. Just-in-time compilation, garbage
collection, address space layout randomisation, CPU frequency scaling, cache
state, background processes and network variability all move the number. A
single run tells you almost nothing.

\begin{paperbox}[Why one run is not a measurement]
Andy Georges, Dries Buytaert and Lieven Eeckhout, \emph{Statistically rigorous
Java performance evaluation}, OOPSLA 2007~\cite{georges2007statistically}.

\medskip
\textbf{Why it matters.} They showed that the then-standard practice --- report
the best of a few runs --- produced conclusions that flatly contradicted
statistically sound analysis of the same data. The paper changed measurement
practice in the systems community, and its argument generalises well beyond
Java.

\medskip
\textbf{What to extract.} The distinction between start-up and steady-state
performance, the need for multiple JVM invocations rather than multiple
iterations within one, and the practice of reporting confidence intervals
rather than a single number.
\end{paperbox}

\begin{center}
\small
\begin{longtable}{@{}L{3.0cm}L{5.2cm}L{6.0cm}@{}}
\toprule
\textbf{Concern} & \textbf{What to do} & \textbf{Why} \\
\midrule
\endfirsthead
\toprule
\textbf{Concern} & \textbf{What to do} & \textbf{Why} \\
\midrule
\endhead
\bottomrule
\endfoot
Warm-up
  & Discard iterations until the metric stabilises; report how you decided
  & Early iterations measure compilation and cache filling, not the system \\
Steady state vs start-up
  & Decide which you are measuring, and say so
  & They are different questions. A serverless workload cares about start-up;
    a long-running service does not \\
Repetitions
  & Many iterations \emph{within} a process, and many process invocations
  & Within-process variation and between-process variation have different
    sources; only the latter captures layout and JIT-decision randomness \\
Interleaving
  & Randomise or alternate the order of configurations
  & Thermal drift and background load change over a session and would otherwise
    confound whichever configuration ran last \\
Environment control
  & Pin frequency, disable turbo, isolate cores, record kernel and library
    versions
  & Otherwise the machine is a nuisance variable you have not controlled \\
Outliers
  & Decide the rule in advance and report how many were removed
  & Removing slow runs after the fact biases towards your preferred result \\
\end{longtable}
\end{center}

\section{Summarising: the Mean That Is Usually Wrong}

\begin{alertbox}[Never take the arithmetic mean of speedups]
This is the most common statistical error in systems papers, and it is easy to
demonstrate.

\medskip
Two benchmarks. On the first, your system takes 1 second where the baseline
takes 2 --- a speedup of 2.0. On the second, yours takes 2 seconds where the
baseline takes 1 --- a speedup of 0.5. The arithmetic mean of the speedups is
$(2.0 + 0.5)/2 = 1.25$, suggesting a 25\% improvement.

\medskip
But total time is 3 seconds for both systems. There is no improvement at all.
Worse, had you defined speedup the other way round --- baseline over yours ---
the arithmetic mean would give the same 1.25 in favour of the baseline. A
summary statistic whose answer depends on which system you put on top is not
measuring anything.

\medskip
The geometric mean, $\left(\prod_{i=1}^{n} s_i\right)^{1/n}$, gives
$\sqrt{2.0 \times 0.5} = 1.0$ either way round. For ratios, use it.
\end{alertbox}

\begin{center}
\small
\begin{tabular}{@{}L{3.6cm}L{10.7cm}@{}}
\toprule
\textbf{Quantity} & \textbf{Summary to use} \\
\midrule
Ratios, speedups, normalised scores & Geometric mean \\
Times, when total time is what matters & Arithmetic mean, or the sum \\
Rates (operations per second)       & Harmonic mean of the rates, equivalently
                                      the reciprocal of mean time \\
Latency distributions               & Median and high percentiles --- p95, p99,
                                      p99.9. The mean of a latency distribution
                                      is close to meaningless, because the tail
                                      is what users experience \\
\bottomrule
\end{tabular}
\end{center}

\begin{conceptbox}[Report an interval, and the whole picture changes]
A table of point values invites the reader to treat a 3\% difference as real. A
table with confidence intervals lets them see that the intervals overlap
entirely, and that the honest conclusion is ``no detectable difference on this
workload''.

\medskip
This is the same discipline as \chref{inference}, applied to systems
measurement, and it has the same consequence: many published performance
improvements would not survive it. Publish the raw per-run measurements
alongside the paper so that others can check --- it costs nothing and it is
what \chref{reproducibility} asks of you.
\end{conceptbox}

\section{Simulation}

When the real system cannot be built, cannot be perturbed, or cannot be run at
the scale of interest, simulation substitutes a model. The model's credibility
is then the whole question.

\begin{center}
\small
\begin{tabular}{@{}L{3.0cm}L{11.3cm}@{}}
\toprule
\textbf{Verification} & Did I build the model \emph{right}? Does the code
                        implement the intended model? Debugging, unit tests,
                        conservation checks, degenerate cases with known
                        answers \\
\textbf{Validation}   & Did I build the \emph{right} model? Does it correspond
                        to the real system well enough for the question being
                        asked? Comparison with measurements from the real
                        system, expert review, sensitivity analysis \\
\bottomrule
\end{tabular}
\end{center}

\begin{alertbox}[Validation is not optional, and it is what gets skipped]
A verified but unvalidated simulation is a correct implementation of an
arbitrary model. Its output is a property of your assumptions, not of the
world, and any conclusion drawn from it is conditional on assumptions nobody
checked.

\medskip
At minimum: compare the model's output against measurements from a real system
in at least one regime where both can be run, and report the discrepancy. Then
run a sensitivity analysis --- vary each assumed parameter and report which
ones the conclusion depends on. If the conclusion flips when an unmeasured
parameter moves within its plausible range, you have found the study's real
limitation, and it belongs in the abstract.
\end{alertbox}

\begin{center}
\small
\begin{tabular}{@{}L{3.4cm}L{10.9cm}@{}}
\toprule
Random seeds       & Fix and publish them; run enough independent replications
                     to estimate variability across seeds \\
Transient removal  & Discard the initial period before the simulated system
                     reaches steady state; state the method used to decide \\
Replications       & Independent runs with different seeds. Report the number
                     and the resulting interval \\
Run length         & Long enough that the estimate stabilises; show that it does \\
\bottomrule
\end{tabular}
\end{center}

\section{Benchmark Suites, and Their Failure Mode}

Standard suites --- SPEC, TPC, YCSB, MLPerf and others --- give comparability
and save enormous effort. They also create an incentive problem.

\begin{conceptbox}[Benchmark overfitting]
When a suite becomes the measure of merit, systems get optimised for the suite
rather than for the workloads it was meant to represent. Over time the suite
stops being a proxy for reality and becomes a target in its own right --- and
improvements on it stop transferring.

\medskip
This is the same phenomenon as test-set overfitting in machine learning
(\chref{mleval}), and it has the same remedies: hold out workloads the
community has not tuned against, refresh the suite periodically, and report
results on a workload of your own alongside the standard one. If your system
wins on the suite and loses on a realistic workload you constructed, say so ---
that is a genuine finding about the suite.
\end{conceptbox}

\begin{traditions}{performance evaluation}
\tradEMP{A benchmark run is an experiment: configurations are treatments, runs
are units, and the analysis needs intervals and repetitions like any other.}
\tradDSR{Benchmarking is summative artificial evaluation of an artefact. The
FEDS caution applies: strong on technical efficacy, weak on whether anyone
benefits in practice.}
\tradQUAL{Asks what practitioners actually optimise for, which is frequently
not what benchmarks measure --- tail latency under partial failure, or
predictability rather than throughput.}
\tradFORM{Provides the asymptotic analysis that measurement complements.
Neither substitutes for the other: an $O(n \log n)$ algorithm can lose to an
$O(n^2)$ one at every $n$ anyone runs.}
\end{traditions}

\begin{lab}[Measure something properly]
Take any two implementations of the same task --- two sorting libraries, two
JSON parsers, two database drivers.

\begin{tightnum}
  \item Write the measurement protocol \emph{first}: warm-up rule, iterations,
    process invocations, interleaving, environment controls, outlier rule.
  \item Run it. Record every individual measurement, not summaries.
  \item Compute the mean, median, and a 95\% confidence interval per
    configuration. Plot the distributions, not just the summaries.
  \item Now do the wrong thing deliberately: report the single best run of
    each. How different is the conclusion?
  \item Add a workload where you expect the winner to lose. Report the
    geometric mean across both.
\end{tightnum}

\medskip
Step 4 is the point of the exercise, and it is worth doing once so that you
never again believe a paper that reports one number.
\end{lab}

\begin{checkpoint}
\begin{tightnum}
  \item Why is the arithmetic mean of speedups wrong, and what test exposes it
    immediately?
  \item Distinguish verification from validation, and say which one an
    unvalidated simulation still passes.
  \item Why are multiple process invocations necessary rather than multiple
    iterations within one process?
  \item Your system wins on SPEC and loses on a production trace. What do you
    report, and what have you actually discovered?
\end{tightnum}
\end{checkpoint}

\begin{chaptersummary}
\begin{tight}
  \item A benchmark is an argument about representativeness. Fix the workload
    set before building, from external evidence, and include a case you expect
    to lose.
  \item One run is not a measurement. Handle warm-up, distinguish start-up from
    steady state, repeat within and across processes, interleave
    configurations, and control the environment.
  \item Geometric mean for ratios; arithmetic for times; harmonic for rates;
    percentiles for latency. The arithmetic mean of speedups is not a
    statistic.
  \item Report intervals, not point values, and publish the raw measurements.
  \item Verification asks whether the model was built right; validation asks
    whether it is the right model. Unvalidated simulation reports your
    assumptions back to you.
  \item Standard suites become targets. Report a workload of your own alongside
    the suite.
\end{tight}
\end{chaptersummary}

\begin{reviewq}
  \item Georges and colleagues showed that a widely used methodology produced
    conclusions contradicting sound analysis of the same data. What does this
    imply about the published performance literature predating it, and what
    should a reviewer today do with such results?
  \item Benchmark suites both enable comparison and distort optimisation.
    Design a governance scheme for a suite in your area that mitigates
    overfitting, and identify who would have to bear its cost.
  \item Simulation validation requires comparison with a real system, but
    simulation is often used precisely because the real system cannot be built.
    How should a study proceed when validation is impossible in principle, and
    what may it then claim?
  \item Tail latency matters more than mean latency to users, yet means
    dominate published results. Give one methodological and one incentive-based
    explanation, and say which you find more persuasive.
\end{reviewq}
